% =====================================================================
%  "Privacy in the Era of Large Opaque Models" workshop variant.
%  Defines everything main.tex pulls in: title, venue string, abstract,
%  introduction, discussion. Nothing else in the paper changes.
% =====================================================================

\newcommand{\wkTitle}{Privacy and Causal Fairness Are Nearly Orthogonal:
Graph Surgery Across Three Generator Families}

\newcommand{\wkVenue}{Privacy in the Era of Large Opaque Models}

% ---------------------------------------------------------------- abstract
\newcommand{\wkAbstract}{%
A differentially private synthetic release is published so that nobody has to look
at the real records again, which makes the release the only thing anyone can
inspect, and makes what it fails to report a problem of its own. We study a
fairness intervention inside that setting. DECAF achieves fair synthetic data by
factoring a GAN's generator along a causal graph and severing the edges that carry
disparity from a protected attribute to an outcome; it has no privacy mechanism,
and the technique has never been shown to work anywhere but in DECAF. We port all three of its fairness definitions to nine generators across three
families (marginals-based, GAN-based and diffusion-based, each with differentially
private variants) over 2{,}520 matched-pair runs at $\varepsilon \in
\{1,10,1000\}$. Three results. The definitions transfer to every family, though the
potency does not: the cut removes $89\%$ of a causal diffusion backbone's disparity
and $29\%$ of a marginals synthesizer's, because a private synthesizer has no
supplied DAG to cut and must bind a structure that noise chose. Privacy and fairness
behave as near-independent knobs: the parity-gap reduction at
$\varepsilon{=}1$ is at least as large as at $\varepsilon{=}1000$ in eight of
eleven generator-dataset cells, and in marginals-based methods the cut costs no
privacy budget at all, because a severed edge is never measured. And DECAF's own argument that privacy follows for free by
post-processing turns out to be sound accounting attached to an unusable output:
the drop-in private discriminator leaves the released table with no signal at all
(lift over the majority class $\le 0.001$) in $49\%$ of runs, a failure two
unrelated changes each repair completely.%
}

% ------------------------------------------------------------------- intro
\newcommand{\wkIntro}{%
A differentially private synthetic release is an opaque object by construction.
The point of publishing it is that the underlying records stay unseen, so every
downstream question, including whether the data is useful and whether a model
trained on it is biased, must be answered against the release itself or not at
all. Two properties of that object therefore matter: what the privacy budget
costs, and whether the release tells the truth about itself.

Fairness sits awkwardly inside this picture. The standard intuition is a
trade-off: DP noise falls hardest on the smallest groups, so tightening
$\varepsilon$ should widen disparities and blunt any attempt to correct them.
Meanwhile the most attractive fairness mechanism for generated data has no privacy
story at all. DECAF \citep{vanbreugel2021decaf} factors a GAN's generator along a
causal DAG, so severing an edge at sampling time makes the generator structurally
unable to use it, and three fairness definitions become three sets of cuts. It is
a single GAN with no privacy guarantee, and nothing about the mechanism has been
demonstrated outside it.

If causal fairness is available only from a generator that cannot offer
differential privacy, the two goals are not merely in tension but architecturally
incompatible. PreFair \citep{pujol2023prefair} shows the incompatibility is not
fundamental, embedding one coarse fairness condition into the private structure
search of one synthesizer, but stops there: one definition, one algorithm, no
evaluation on any other.

We ask whether DECAF's edge surgery is a property of DECAF or of causal
factorisation, and what a privacy budget does to it, attaching all three
definitions to nine generators from three unrelated families. That includes
marginals-based synthesizers, where a severed edge is never privately measured in
the first place.%
}

% -------------------------------------------------------------- discussion
\newcommand{\wkDiscussion}{%
\textbf{Post-processing is an accounting argument, not a guarantee that anything
comes out.} DECAF states that it ``can be directly used for generating private
synthetic data by replacing the standard discriminator by a differentially private
discriminator,'' citing the post-processing theorem. The privacy claim is
correct and we do not dispute it. What our runs show is that the sentence is
nonetheless the wrong advice: done literally, it produces a table carrying no
downstream signal in $49\%$ of runs across the two DP-GAN backbones, rising to
$80\%$ for the DP-CTGAN backbone on COMPAS at $\varepsilon{=}1$. The failure
is narrow and locatable. Keeping adversarial training but changing the data
representation fixes it, and keeping the representation but dropping adversarial
training for diffusion fixes it; two unrelated repairs of one symptom put the
cause in DP-SGD applied to a weight-clipped critic. The general lesson is worth
stating plainly, because this style of argument is common: post-processing
composes privacy budgets, and says nothing whatever about whether the composed
system still generates data. Papers that reach for it owe an experiment.

\textbf{The fairness constraint is free in the accounting, but not in fidelity.}
In the marginals-based family a severed edge is excluded from candidacy before the
private structure search runs, so no $\varepsilon$ is ever spent measuring it; in
the DECAF-style families the cut happens after training and is itself
post-processing. Both are free in the privacy accounting, for different reasons,
and against the standard picture in which fairness is a tax paid out of the privacy
budget that is a strong negative result. It is worth being precise about what the
freed budget does \emph{not} do, though: it does not buy a better-fitting release.
Across the marginals-based generators, matched-pair 2-way TVD is slightly worse
under the cut, by $0.001$ for FTU, $0.003$ for CF and $0.008$ for DemPar --- ordered
by how much structure each removes.
The reason is structural rather than accounting: MST fixes its per-round
$\varepsilon$ from the number of edges a spanning tree needs, before selection
begins, so excluding candidates leaves no extra budget for the edges that remain
and merely shrinks the pool each round draws from. Prune hard enough and a round
has no admissible candidate at all, so the model ends as a spanning forest rather
than a tree --- less structure for the same $\varepsilon$. Fairness costs no
privacy here; it still costs a little fidelity.

\textbf{Opacity has a measurable direction, and it flatters.} DECAF's Definition 4
warns that a distribution can look fair scored against itself and still yield
disparate impact under the distribution the model is deployed into. We can put a
number on it. Every fairness result above is measured against real held-out
people; an analyst holding only the release must score the model against the
synthetic data instead. Under DemPar on Adult, the release reports a parity gap of
$0.010$ while the same models sit at $0.035$ against real people, a $3.5\times$
understatement. The direction is the uncomfortable part: the distortion is
\emph{largest} for the mechanism that most reduces the true gap, because the
intervention acts on the generative process and so is fully realised in the
synthetic distribution, while the deployed model meets a world still containing
every proxy the generator suppressed. A private release used as its own audit
substrate will certify models it should not, and the better the fairness
mechanism, the more confidently it will do so. Reporting both references costs one
extra evaluation on data the publisher already holds.

\textbf{A tight privacy budget manufactures fairness that is not there.} On Adult
at $\varepsilon{=}10$, both private GAN backbones report a parity gap of $0.000$
before any mechanism is applied, with lift at or below $0.001$ over the majority
class. A table that teaches a classifier nothing contains no disparity to find.
Any privacy-fairness benchmark that reports gaps without utility beside them is
unfalsifiable in exactly the regime it exists to study.%
}
